\subsection{Rancangan Segmentasi Data \textit{Off-Chain}}
\label{subsec:rancangan-segmentasi-offchain}

% Kontrol akses granular tidak dapat ditegakkan hanya dengan mengganti token apabila data RME masih disimpan sebagai satu \textit{encrypted blob} yang terlalu besar. Oleh karena itu, rancangan ini memecah data \textit{off-chain} menjadi beberapa unit data terenkripsi berdasarkan kategori layanan dan kebutuhan akses. Segmentasi tersebut mengikuti pembagian isi RME yang telah dibahas pada Subbab \ref{subsec:isi-rme}, yaitu dokumentasi administratif dan dokumentasi klinis, serta mengacu pada kelompok data seperti rawat jalan, laboratorium, dan apotek.

% Setiap segmen data di IPFS memiliki CID dan metadata sendiri di IOTA. Metadata tersebut memuat informasi seperti \texttt{dataset\_category}, \texttt{ipfs\_cid}, \texttt{integrity\_hash}, serta kapsul kunci yang diperlukan oleh PRE. Dengan demikian, ketika seorang aktor meminta akses terhadap suatu segmen, \textit{server} PRE dapat mencocokkan \textit{caveats} pada token dengan metadata segmen data yang diminta. Apabila token hanya mengizinkan akses terhadap data laboratorium, maka token tersebut tidak dapat digunakan untuk meminta segmen riwayat klinis rawat jalan yang tidak relevan.

% Pada skenario layanan kesehatan kolaboratif, segmentasi data ini menghasilkan pemisahan akses sebagai berikut.

% \begin{enumerate}[leftmargin=*]
% 	\item Petugas administrasi mengakses segmen data administratif dan pendaftaran kunjungan.
% 	\item Dokter penanggung jawab pasien mengakses segmen klinis yang relevan, seperti rawat jalan dan hasil pemeriksaan penunjang yang diperlukan untuk pengambilan keputusan klinis.
% 	\item Petugas laboratorium mengakses segmen permintaan pemeriksaan dan menulis segmen hasil pemeriksaan tanpa membuka keseluruhan riwayat klinis pasien.
% \end{enumerate}

% Rancangan ini menjadikan segmentasi penyimpanan sebagai prasyarat teknis bagi penegakan prinsip \textit{need-to-know}. Skema penyimpanan rinci dapat dilihat pada Lampiran \ref{lampiran:skema-penyimpanan}.
